% =====================================================================
% Revised model and method sections. Horizon is a single 48-hour block,
% so the day index is dropped and intervals are indexed by t alone.
% Carbon limits are nodal.
%
% Replaces: Sec. II-B-3 boundary conditions and Sec. II-B-6.
% Adds:     Sec. III.
% New symbols for the nomenclature are listed at the end.
% =====================================================================

\section{Carbon-Constrained Shared-Storage Planning}
\label{sec:model}

% ---------------------------------------------------------------------
\subsection{Horizon and Carbon Vintages}
\label{sec:horizon}

Each scenario $s\in\mathcal{S}$ is a $48$-hour block,
$t\in\mathcal{T}=\{1,\dots,T\}$ with $T=48$, weighted by an annual
occurrence factor $\rho_s=8760/(T\Delta t)$ so that operating costs
annualize consistently. Two days are long enough to contain a full
charge--discharge cycle together with the overnight interval in which
stored energy is most valuable, and short enough that a vintage-resolved
inventory remains tractable, since the number of vintage layers grows
with the horizon.

The stored energy at each candidate bus is partitioned into vintages
indexed by $k\in\mathcal{K}=\{0,1,\dots,T\}$. Layer $k=0$ carries the
inventory present at the start of the block; layer $k\geq1$ is created by
charging in interval $k$ and, once created, retains the carbon intensity
$\mathscr{w}_{i,k,s}^{\mathrm{es}}$ of the bus at that interval until it
is discharged. The layer-wise energy and carbon dynamics
\eqref{eq:layer_energy_dynamic}--\eqref{eq:layer_carbon_dynamic} and the
activation logic \eqref{eq:layer_activation_lower}--%
\eqref{eq:carbon_intensity_bound} are unchanged; with a single block the
inter-day transfer constraints are no longer required, and the cyclic
conditions apply at the ends of the block,
\begin{equation}
\sum_{k\in\mathcal{K}}e_{i,k,s,T+1}^{\mathrm{es}}
=e_{i,0,s,1}^{\mathrm{es}},
\qquad
\sum_{k\in\mathcal{K}}m_{i,k,s,T+1}^{\mathrm{es}}
=m_{i,0,s,1}^{\mathrm{es}}.
\label{eq:block_cycle}
\end{equation}

% ---------------------------------------------------------------------
\subsection{Nodal Carbon Flow}
\label{sec:nodal_carbon}

Carbon is traced through the network by proportional sharing
\cite{chen2025copf}. The carbon injected at bus $i$ by storage discharge
carries the intensity of the vintage supplying it,
\begin{equation}
R_{i,s,t}^{\mathrm{es,dc}}
=
\sum_{k\in\mathcal{K}}
\mathscr{w}_{i,k,s}^{\mathrm{es}}
P_{i,k,s,t}^{\mathrm{es,dc}},
\label{eq:storage_carbon_flow}
\end{equation}
which is what distinguishes this model from a well-mixed inventory: the
carbon released depends on \emph{which} stored energy is used, not only
on how much.

Collecting the power and carbon arriving at bus $i$ from every source,
\begin{align}
A_{i,s,t}
=\;&
\sum_{g\in\mathcal{G}_i}P_{g,s,t}
+P_{i,s,t}^{\mathrm{PV}}
+\delta_{i0}P_{s,t}^{\mathrm{grid}}
+P_{i,s,t}^{\mathrm{es,dc}}
\nonumber\\
&+\sum_{h:(h,i)\in\mathcal{L}}P_{hi,s,t}^{+}
+\sum_{j:(i,j)\in\mathcal{L}}P_{ij,s,t}^{-},
\label{eq:incoming_power}\\
B_{i,s,t}
=\;&
\sum_{g\in\mathcal{G}_i}\mathscr{w}_{g}^{G}P_{g,s,t}
+\delta_{i0}\mathscr{w}_{s,t}^{\mathrm{grid}}P_{s,t}^{\mathrm{grid}}
+R_{i,s,t}^{\mathrm{es,dc}}
\nonumber\\
&+\sum_{h:(h,i)\in\mathcal{L}}\mathscr{w}_{h,s,t}P_{hi,s,t}^{+}
+\sum_{j:(i,j)\in\mathcal{L}}\mathscr{w}_{j,s,t}P_{ij,s,t}^{-},
\label{eq:incoming_carbon}
\end{align}
the nodal intensity is defined by the mixing identity
\begin{equation}
\mathscr{w}_{i,s,t}A_{i,s,t}=B_{i,s,t}.
\label{eq:nodal_mixing}
\end{equation}
Photovoltaic output appears in $A_{i,s,t}$ but not in $B_{i,s,t}$ and so
dilutes the nodal intensity, which is the mechanism by which local
generation relaxes a nodal limit.

% ---------------------------------------------------------------------
\subsection{Nodal Carbon Targets}
\label{sec:carbon_target}

The data-center bus is held to a target $\overline{\mathscr{w}}$, and the
remaining buses to $\overline{\mathscr{w}}^{\mathrm{o}}$. Imposing these
as hard limits is unsatisfactory for a planning study. A target loose
enough to be always attainable does not influence the design, and one
tight enough to influence it is, on this system, already tight enough to
be unreachable in some intervals---at which point the planning problem
returns no design at all rather than an expensive one, and the
sensitivity of the design to the target cannot be traced.

We therefore admit a nodal excess $\Xi_{i,s,t}\geq0$, in the units of
$B_{i,s,t}$, and price it:
\begin{equation}
\mathscr{w}_{i,s,t}A_{i,s,t}
\leq
\overline{\mathscr{w}}_i A_{i,s,t}
+
\Xi_{i,s,t},
\label{eq:nodal_carbon_target}
\end{equation}
with $\overline{\mathscr{w}}_i=\overline{\mathscr{w}}$ at
$i=i^{\mathrm{DC}}$ and $\overline{\mathscr{w}}^{\mathrm{o}}$ elsewhere.
The interval cost \eqref{eq:interval_cost} gains
$\sum_{i\in\mathcal{N}}c^{\mathrm{CO_2}}\Xi_{i,s,t}$.

Pricing the excess in the units of the quantity it measures is deliberate
and is not equivalent to the customary big-$M$ relaxation, in which a
dimensionless violation variable is scaled by a bound on nodal power
before entering the constraint. Under that form the implied price per
unit of carbon is inversely proportional to a numerical bound, so
re-deriving the bound changes the cost of a violation without changing
its physical size, and the value attributed to storage moves with it.
Here $c^{\mathrm{CO_2}}$ is an assumption of the study, stated once. The
formulation retains both limits: $c^{\mathrm{CO_2}}\!\to\!\infty$ recovers
the hard cap wherever it is attainable, and $c^{\mathrm{CO_2}}\!=\!0$
recovers carbon-agnostic planning, so the influence of the target on
siting and sizing can be traced continuously between them.

% ---------------------------------------------------------------------
\subsection{Where the Storage Value Originates}
\label{sec:value_origin}

Substituting \eqref{eq:nodal_mixing} and \eqref{eq:storage_carbon_flow}
into \eqref{eq:nodal_carbon_target} and separating the terms that do not
involve storage discharge isolates what storage contributes. With
$A_{i,s,t}^{\setminus}$ and $B_{i,s,t}^{\setminus}$ denoting
\eqref{eq:incoming_power} and \eqref{eq:incoming_carbon} without their
storage-discharge terms, and
\begin{equation}
D_{i,s,t}
=
B_{i,s,t}^{\setminus}
-
\overline{\mathscr{w}}_i A_{i,s,t}^{\setminus},
\label{eq:carbon_deficit}
\end{equation}
the target reads
\begin{equation}
\sum_{k\in\mathcal{K}}
\bigl(\overline{\mathscr{w}}_i-\mathscr{w}_{i,k,s}^{\mathrm{es}}\bigr)
P_{i,k,s,t}^{\mathrm{es,dc}}
+
\Xi_{i,s,t}
\geq
D_{i,s,t}.
\label{eq:carbon_gap_priced}
\end{equation}

Equation \eqref{eq:carbon_gap_priced} is the mechanism the model exists
to represent. A positive $D_{i,s,t}$ is a compliance deficit that must be
met either by discharging a vintage whose intensity lies below the
target---at a rate weighted by how far below---or by paying for the
excess. The planner is therefore rewarded for charging during
low-intensity intervals specifically so that the resulting vintages are
available when $D_{i,s,t}$ is positive, and the worth of an installation
depends on the intensities of the vintages it can hold, not on its energy
capacity alone. A well-mixed inventory assigns one intensity to all
stored energy and cannot distinguish a battery charged at the cleanest
hours from one charged at the dirtiest; the coefficient
$\overline{\mathscr{w}}_i-\mathscr{w}_{i,k,s}^{\mathrm{es}}$ is exactly
what that abstraction discards.

% ---------------------------------------------------------------------
\subsection{Tractability of the Nodal Identity}
\label{sec:tractability}

Equations \eqref{eq:nodal_mixing} and \eqref{eq:carbon_energy_relation}
are bilinear, so the planning problem is a mixed-integer nonlinear
program. Two treatments are available and they are not interchangeable.

Solving the bilinear equalities exactly requires a nonconvex
mixed-integer solver. This is the reference formulation and we use it to
verify the relaxation below, but its cost grows steeply once the target
binds: on instances where a relaxed solve terminates in under a minute,
the exact solve does not converge within an hour.

The customary alternative replaces each bilinear product by its
McCormick envelope. Care is required, because the envelope's usefulness
here is one-sided. Enforcing \eqref{eq:nodal_carbon_target} needs a lower
bound on $\mathscr{w}_{i,s,t}$, which the envelope supplies as
$\mathscr{w}_{i,s,t}\geq B_{i,s,t}/\overline{A}_i$ where
$\overline{A}_i$ bounds the incoming power. That bound is weak in
proportion to the slack in $\overline{A}_i$, and its companion lower
bound on the product,
$B\geq\overline{\mathscr{w}}A+\overline{A}\mathscr{w}
-\overline{\mathscr{w}}\overline{A}$, becomes vacuous whenever $A$ sits
well below $\overline{A}$. A bound taken from line thermal ratings, which
exceed realized flows by an order of magnitude on a distribution feeder,
therefore admits solutions in which a line carries no carbon at all and
the nodal target is satisfied at any level.

Bounds must instead be derived from what a bus can absorb. Rearranging
the nodal power balance gives the identity
\begin{equation}
A_{i,s,t}
=
P_{i,s,t}^{L}
+P_{i,s,t}^{\mathrm{es,ch}}
+\!\!\sum_{h:(h,i)\in\mathcal{L}}\!\!P_{hi,s,t}^{-}
+\!\!\sum_{j:(i,j)\in\mathcal{L}}\!\!P_{ij,s,t}^{+},
\label{eq:incoming_identity}
\end{equation}
whose terms are individually bounded by the bus load, the installed
charging power, and the demand and injection of the subtree below each
line. Loads are data rather than variables, so the first term is exact.
The resulting $\overline{A}_i$ is tighter than a thermal rating by roughly
an order of magnitude, and because it is combined with the rating by
minimum it can only tighten the relaxation, never admit a solution the
rating excluded.

The tightness of $\overline{A}_i$ should be reported alongside any result
obtained from the relaxed model, and the relaxation should be checked
against the exact formulation on reduced instances. A carbon target that
does not alter the objective as it is tightened is not evidence that
storage is unnecessary; it is evidence that the relaxation has not been
tightened enough to enforce it.

% =====================================================================
\section{Decision-Focused Scenario Selection}
\label{sec:dfl}

% ---------------------------------------------------------------------
\subsection{The Selection Problem}
\label{sec:selection_problem}

Let $\xi$ denote an operating scenario---a joint realization over the
48-hour block of load, PV availability, workload arrivals, grid carbon
intensity and prices---distributed as $\mathbb{P}$. Writing
$c(\boldsymbol{x}^{\mathrm{pl}},\xi)$ for the optimal operating cost
under design $\boldsymbol{x}^{\mathrm{pl}}$ and scenario $\xi$, the
planning problem of interest is
\begin{equation}
\min_{\boldsymbol{x}^{\mathrm{pl}}}
\quad
C^{\mathrm{inv}}(\boldsymbol{x}^{\mathrm{pl}})
+
\mathbb{E}_{\xi\sim\mathbb{P}}
\bigl[c(\boldsymbol{x}^{\mathrm{pl}},\xi)\bigr].
\label{eq:true_problem}
\end{equation}
Each evaluation of $c$ solves a mixed-integer program over the whole
block, so \eqref{eq:true_problem} is solved in practice on a support of
$K$ scenarios $\mathcal{X}=\{\xi_1,\dots,\xi_K\}$ with weights
$\boldsymbol{\pi}\in\Delta^{K-1}$:
\begin{equation}
\boldsymbol{x}^{\star}(\mathcal{X},\boldsymbol{\pi})
\in
\argmin_{\boldsymbol{x}^{\mathrm{pl}}}
\;
C^{\mathrm{inv}}(\boldsymbol{x}^{\mathrm{pl}})
+
\sum_{k=1}^{K}\pi_k\,c(\boldsymbol{x}^{\mathrm{pl}},\xi_k).
\label{eq:support_problem}
\end{equation}

The support is conventionally chosen to represent $\mathbb{P}$ in
distribution, by clustering or by maximizing coverage. That criterion is
indifferent to what the support is for. Two supports equally
representative in a distributional sense can induce designs of very
different quality, because \eqref{eq:support_problem} responds only to
the features of $\xi$ that bind its constraints---and under
\eqref{eq:carbon_gap_priced} the binding feature is the coincidence of a
compliance deficit with the availability of low-intensity vintages, which
no distributional distance measures. We therefore select the support by
the quality of the decision it induces:
\begin{equation}
\min_{\mathcal{X},\boldsymbol{\pi}}
\quad
\mathcal{L}(\mathcal{X},\boldsymbol{\pi})
:=
\mathbb{E}_{\xi\sim\mathbb{P}}
\bigl[
c\bigl(\boldsymbol{x}^{\star}(\mathcal{X},\boldsymbol{\pi}),\xi\bigr)
\bigr].
\label{eq:decision_loss}
\end{equation}
The support enters \eqref{eq:decision_loss} only through the design it
produces; a scenario that is unusual but does not change
$\boldsymbol{x}^{\star}$ has no effect on the loss.

% ---------------------------------------------------------------------
\subsection{Generative Parameterization}
\label{sec:parameterization}

Searching over $\mathcal{X}$ directly is unattractive: scenarios are
high-dimensional and must remain physically admissible, and restricting
the search to observed scenarios reduces \eqref{eq:decision_loss} to a
combinatorial subset selection over the historical record.

We search instead in the latent space of a conditional generator
$G_{\phi}$, trained beforehand on the historical pool, which maps a
latent $\boldsymbol{z}\in\mathbb{R}^{n_z}$ and a context
$\boldsymbol{\omega}$ to a scenario. With
$\boldsymbol{Z}=[\boldsymbol{z}_1,\dots,\boldsymbol{z}_K]$ and
$\boldsymbol{\Omega}=[\boldsymbol{\omega}_1,\dots,\boldsymbol{\omega}_K]$,
\begin{equation}
\mathcal{X}(\boldsymbol{Z},\boldsymbol{\Omega})
=
\bigl\{
G_{\phi}(\boldsymbol{z}_k,\boldsymbol{\omega}_k)
\bigr\}_{k=1}^{K},
\label{eq:generated_support}
\end{equation}
and $\boldsymbol{\pi}=\mathrm{softmax}(\boldsymbol{\lambda})$. The
generator is held fixed; only $(\boldsymbol{Z},\boldsymbol{\lambda})$ are
optimized. The support may thus leave the historical record while
remaining on the manifold the generator learned, and the search is over
$Kn_z$ continuous variables rather than over subsets.

The contexts $\boldsymbol{\Omega}$ are taken from $K$ observed scenarios
selected by a coverage rule and then held fixed, anchoring each support
member to a distinct region of the conditioning space. This is a
restriction and should be stated as one: any attribute of $\xi$ carried
by the context rather than by the latent cannot be altered by the search,
so the reachable set is a neighbourhood of the anchors rather than the
generator's full range. Whether that neighbourhood contains supports of
the quality the method needs is a property of the generator and the
anchors, not of the optimizer, and is checked directly by comparing the
decision-relevant statistics of generated scenarios against observed ones.

% ---------------------------------------------------------------------
\subsection{Score-Function Optimization}
\label{sec:score_function}

The composition in \eqref{eq:decision_loss} is not differentiable in
$(\boldsymbol{Z},\boldsymbol{\lambda})$. The map from support to design
is the $\argmin$ of a mixed-integer program and is piecewise constant:
almost everywhere its gradient vanishes, and where the active set changes
it does not exist. Neither differentiating through the solver nor
smoothing its optimality conditions applies once integer siting variables
are present.

We therefore treat the support as a random variable and optimize the
expected loss. With
\begin{equation}
q_{\boldsymbol{\theta}}(\boldsymbol{Z})
=
\mathcal{N}\bigl(\boldsymbol{Z};\boldsymbol{\mu},\sigma^2\boldsymbol{I}\bigr),
\qquad
\boldsymbol{\theta}=(\boldsymbol{\mu},\boldsymbol{\lambda}),
\label{eq:policy}
\end{equation}
and $\mathcal{J}(\boldsymbol{\theta})=
\mathbb{E}_{\boldsymbol{Z}\sim q_{\boldsymbol{\theta}}}[\mathcal{L}]$,
the score-function identity gives
\begin{equation}
\nabla_{\boldsymbol{\theta}}\mathcal{J}
=
\mathbb{E}_{\boldsymbol{Z}\sim q_{\boldsymbol{\theta}}}
\bigl[
\mathcal{L}\,
\nabla_{\boldsymbol{\theta}}\log q_{\boldsymbol{\theta}}(\boldsymbol{Z})
\bigr],
\label{eq:score_function}
\end{equation}
which requires only evaluations of $\mathcal{L}$, never its derivative.
With $N$ samples per iteration and a baseline $b$ tracking the running
mean of the loss,
\begin{equation}
\widehat{\nabla_{\boldsymbol{\theta}}\mathcal{J}}
=
\frac{1}{N}\sum_{n=1}^{N}
\bigl(\mathcal{L}^{(n)}-b\bigr)
\nabla_{\boldsymbol{\theta}}
\log q_{\boldsymbol{\theta}}(\boldsymbol{Z}^{(n)}).
\label{eq:reinforce}
\end{equation}
The baseline leaves the estimator unbiased and reduces its variance,
which matters because each $\mathcal{L}^{(n)}$ costs one solve of
\eqref{eq:support_problem} followed by a fixed-design evaluation, so $N$
is necessarily small. The expectation in \eqref{eq:decision_loss} is
replaced during training by an average over a fixed validation set
$\mathcal{V}$, held constant across iterations so that losses from
different iterations are comparable, and disjoint from the test set on
which results are reported.

% ---------------------------------------------------------------------
\subsection{Resolution Limits}
\label{sec:resolution}

Two properties of \eqref{eq:reinforce} determine whether the procedure
can work at all, independently of how it is tuned, and both follow from
$\mathcal{L}$ being the output of a terminated mixed-integer solve.

Each $\mathcal{L}^{(n)}$ is obtained at a relative optimality gap
$\epsilon$. Differences between samples smaller than
$\epsilon\mathcal{L}$ are therefore ordered not by the decision but by
where the solver stopped, and if the spread of $\mathcal{L}^{(n)}$ within
an iteration is comparable to $\epsilon\mathcal{L}$, \eqref{eq:reinforce}
estimates the gradient of numerical noise. The gap must be chosen against
the spread of the losses being compared rather than against the magnitude
of the objective; the two differ by orders of magnitude here, because the
storage decision moves a small fraction of total system cost while every
design pays the same untouchable remainder.

The same threshold governs convergence. An update changing the loss by
less than $\epsilon\mathcal{L}$ is indistinguishable from no update, so a
run whose total improvement falls below that threshold has not been shown
to have learned anything, however smooth its trajectory appears. We
therefore report improvement relative to $\epsilon\mathcal{L}$ rather
than in absolute terms, and report the displacement
$\|\boldsymbol{\mu}-\boldsymbol{\mu}_0\|$ of the policy mean alongside
it: a run in which the loss did not move while the mean did was
searching, whereas one whose mean travelled less than the width $\sigma$
it samples at never left its starting neighbourhood.

% ---------------------------------------------------------------------
\subsection{Selection Among Candidates}
\label{sec:finalists}

A stochastic policy visits many supports and the last iterate need not be
the best. We retain the $F$ distinct designs of lowest validation loss
encountered during training, together with the design induced by the
deterministic support $\boldsymbol{\mu}$, and select among them on
$\mathcal{V}$. Candidates differing by less than $\epsilon\mathcal{L}$
are treated as tied, and ties are broken toward the smaller installation,
which avoids reporting capacity the evaluation cannot show to be
worthwhile.

This step is a selection over sampled candidates, not a consequence of
the gradient, and it accounts for part of any advantage over selecting a
support at random. Separating the two contributions requires a control
that draws the same number of candidates from $q_{\boldsymbol{\theta}}$
with $\boldsymbol{\theta}$ held at its initial value and applies the same
selection: what that control does not achieve is what the policy update
contributed.

% ---------------------------------------------------------------------
\subsection{Evaluation Protocol}
\label{sec:protocol}

All methods are compared by the value of the design each produces, on one
test subset $\mathcal{T}^{\mathrm{te}}$ and against one reference:
\begin{equation}
V_{\mathrm{m}}
=
\Bigl[
\min_{\boldsymbol{y}^{\mathrm{op}}}
\!\!\sum_{\xi\in\mathcal{T}^{\mathrm{te}}}\!\!
c(\boldsymbol{0},\xi)
\Bigr]
-
\Bigl[
C^{\mathrm{inv}}(\boldsymbol{x}^{\star}_{\mathrm{m}})
+
\!\!\sum_{\xi\in\mathcal{T}^{\mathrm{te}}}\!\!
c(\boldsymbol{x}^{\star}_{\mathrm{m}},\xi)
\Bigr].
\label{eq:storage_value}
\end{equation}
Total system cost is a poor denominator, since the storage decision moves
a few percent of it and every design pays the same remainder;
$V_{\mathrm{m}}$ measures what installing storage is worth, net of its own
investment. The first term is common to all methods, so ranking by
$V_{\mathrm{m}}$ and by out-of-sample cost is identical---and a difference
in $V$ is meaningful only when both designs were evaluated against the
same reference on the same subset, which we verify rather than assume.

Two quantities bound the comparison. The perfect-information value
$V^{\mathrm{PI}}$, obtained by optimizing the design directly on
$\mathcal{T}^{\mathrm{te}}$, is attainable by no selection rule and
converts values into captured fractions; where the solve terminates early
we report the interval implied by its incumbent and dual bound rather
than a point estimate. And because a stochastic method must be compared
against a distribution rather than a draw, we report the spread of $V$
over repeated runs alongside its mean, and take as the reference point
the strongest rule requiring neither learning nor a favourable draw.

% =====================================================================
% Additional nomenclature
% =====================================================================
% \item[$\Xi_{i,s,t}$]      Carbon emitted at bus $i$ beyond its target.
% \item[$c^{\mathrm{CO_2}}$] Price per unit of excess carbon.
% \item[$A_{i,s,t},B_{i,s,t}$] Power and carbon arriving at bus $i$.
% \item[$\overline{A}_i$]   Upper bound on incoming power at bus $i$.
% \item[$D_{i,s,t}$]        Compliance deficit before storage discharge.
% \item[$\mathcal{X},\boldsymbol{\pi}$] Support scenarios and weights.
% \item[$G_{\phi}$]         Pretrained conditional scenario generator.
% \item[$\boldsymbol{Z},\boldsymbol{\Omega}$] Support latents and contexts.
% \item[$q_{\boldsymbol{\theta}}$] Sampling policy over supports.
% \item[$\mathcal{L}$]      Decision loss of a support.
% \item[$\mathcal{V}$]      Fixed validation set used during training.
% \item[$\epsilon$]         Relative optimality gap of the planning solve.
% \item[$V_{\mathrm{m}},V^{\mathrm{PI}}$] Storage value; its
%                           perfect-information bound.
